%% appendix/C-code-listings.tex
%% Appendix C: Selected Code Listings
%% Pulls out the most important code snippets from the eight notebooks for
%% the reader who wants to see the experiments in one place. Modeled on
%% Joseph Genzer's Appendix B (Code).
%%
%% Once the notebooks are written, the four placeholder listings below will
%% be replaced with the real code, formatted via the thesisPython lstlisting
%% style defined in style/thesis.sty.

\section{Selected Code Listings}\label{app:code}

\todo{Include 4 to 6 of the most informative snippets from the notebooks
(not full notebooks). Each listing should be self-contained at the function
level, captioned, and referenced from the main text where the corresponding
figure appears.

Suggested listings to include here:
\begin{itemize}
  \item C.1 One-spike Johnstone simulation
        (driving function for Section~\ref{subsec:bbp}).
        Source: \texttt{notebooks/04\_section4\_demos.py}.
  \item C.2 Eigenvector overlap simulation
        (Section~\ref{subsec:eigenvector}).
        Source: \texttt{notebooks/07\_eigenvector\_overlap.ipynb}.
  \item C.3 Strong-spike fluctuation simulation
        (Section~\ref{subsec:fluctuations}).
        Source: \texttt{notebooks/06\_fluctuation\_qq\_plots.ipynb}.
  \item C.4 Marchenko-Pastur bulk simulation
        (Section~\ref{subsec:scm-highdim}).
        Source: \texttt{notebooks/04\_marchenko\_pastur\_tracy\_widom.ipynb}.
\end{itemize}

Each listing will be wrapped in a \texttt{lstlisting} environment using the
\texttt{thesisPython} style defined in \texttt{style/thesis.sty}, with a
caption and a \texttt{label} for cross-referencing from the main text.}

\clearpage
